\documentclass[11pt,a4paper]{article}
\usepackage[margin=20mm]{geometry}
\usepackage{fontspec}
\setmainfont{DejaVu Serif}
\setsansfont{DejaVu Sans}
\setmonofont{DejaVu Sans Mono}[Scale=0.85]
\usepackage{polyglossia}
\setdefaultlanguage{russian}
\usepackage{graphicx}
\usepackage{hyperref}
\usepackage{enumitem}
\usepackage{fancyvrb}
\hypersetup{hidelinks}
\setlength{\parindent}{0pt}
\setlength{\parskip}{4pt}
\usepackage{titlesec}
\titleformat{\section}{\large\bfseries}{}{0pt}{}
\titlespacing*{\section}{0pt}{10pt}{4pt}
\sloppy
\setlist{nosep,leftmargin=*}
\emergencystretch=3em
\newcommand{\shot}[2]{\begin{center}\includegraphics[width=\linewidth,height=.34\textheight,keepaspectratio]{#1}\par\small #2\end{center}}
\begin{document}
\begin{center}
{\small Университет ИТМО · Фронт-энд разработка · 2026}\par
{\LARGE Домашняя работа № 4}\par
{\large SVG-спрайт}\par
Светлов Ярослав, группа J3211
\end{center}
\section*{Цель}
Собрать набор из 3–5 иконок в общий SVG-спрайт и использовать его в приложении. Работа сохранена как независимая копия ДЗ3.
\section*{Реализация}
В \texttt{public/sprite.svg} объявлены пять символов: пользователь, поиск, модель, файл и выход. Разметка обращается к ним через внешний элемент use, например \texttt{/sprite.svg\#file}. Фигуры не дублируются в HTML.

Цвет наследуется через currentColor, размер задаёт класс .icon. Иконки следуют светлой и тёмной темам. Они декоративные: используются aria-hidden="true" и focusable="false"; рядом остаются текстовые названия действий.
\section*{Проверка}
В Firefox 148 проверены четыре основных экрана в обеих темах, фильтры, сброс и восстановление сессии. Результаты: \texttt{checks/browser.json}.

В Chromium проверено использование всех пяти идентификаторов. На ширине 390 px нет горизонтального переполнения всей страницы, таблицы имеют собственную прокрутку. Девять API-тестов пройдены.
\section*{Запуск и проверка}
Требуется Node.js 22.12 или новее. Команды выполняются из папки этой работы.
\begin{verbatim}
npm ci
npm start
\end{verbatim}
Открыть \url{http://localhost:3004}. Автотесты: \texttt{npm test}.
Учебные аккаунты: \texttt{admin@nexus.ai / password} и \texttt{vision@nexus.ai / cv12345}. Данные сохраняются в \texttt{db.local.json}, исходная база \texttt{db.json} не меняется.\section*{Вывод}
Создан единый спрайт из пяти иконок. Их размер и цвет управляются CSS; иконки не заменяют доступные текстовые подписи.

\clearpage
\section*{Скриншоты}
\shot{checks/light.png}{Рисунок 1. Иконки в светлой теме.}
\shot{checks/dark.png}{Рисунок 2. Иконки в тёмной теме.}

\end{document}
